\documentclass[12pt]{article}

% Ћирилица: кодирање фонта и језик (српски, ћирилица)
\usepackage[T2A]{fontenc}
\usepackage[utf8]{inputenc}
\usepackage[serbianc]{babel}

% Величина стране и маргине
\usepackage[a4paper,top=2cm,bottom=2cm,left=2.5cm,right=2.5cm]{geometry}

% Корисни пакети
\usepackage{amsmath}
\usepackage{enumitem}
\usepackage{booktabs}
\usepackage{tcolorbox}
\tcbuselibrary{breakable}
\usepackage[colorlinks=true, allcolors=blue]{hyperref}

% Цртање електричних кола (европски симболи)
\usepackage[european]{circuitikz}
\usetikzlibrary{babel}

% Оквир за решене примере
\newtcolorbox[auto counter]{primer}{breakable,
    colback=white, colframe=black!55, colbacktitle=black!55, coltitle=white,
    fonttitle=\bfseries\small, title={Пример~\thetcbcounter},
    boxrule=0.5pt, arc=1.5pt, left=6pt, right=6pt, top=3pt, bottom=3pt,
    before skip=8pt, after skip=10pt}

% Оквир за кључне формуле
\newtcolorbox{kljucno}{colback=blue!4, colframe=blue!45!black,
    boxrule=0.5pt, arc=1.5pt, left=8pt, right=8pt, top=4pt, bottom=4pt}

\title{Везе отпорника и кондензатора}
\author{Основи електротехнике 1}
\date{}

\begin{document}
\maketitle

%==================================================================
\section{Редна веза отпорника}
%==================================================================

Отпорници су везани \textbf{редно} када су повезани један за другим, тако да
између тачака $A$ и $B$ постоји \textbf{само један пут} за струју.

\begin{center}
\begin{circuitikz}
    \draw (0,0) node[left]{$A$} to[short, o-] (0.5,0)
          to[R, l=$R_1$] (2.5,0)
          to[R, l=$R_2$] (4.5,0)
          to[R, l=$R_3$] (6.5,0)
          to[short, -o] (7,0) node[right]{$B$};
\end{circuitikz}
\end{center}

Пошто струја нема где да скрене, \textbf{кроз све отпорнике тече иста струја}.
Сваки отпорник отежава пролаз струје, па се отпорности сабирају:

\begin{kljucno}
\[
R_e = R_1 + R_2 + \dots + R_n
\]
\end{kljucno}

$R_e$ је \textbf{еквивалентна отпорност}. То је један отпорник којим можемо
заменити све редно везане отпорнике, а да се у остатку кола ништа не промени.

\begin{itemize}
    \item Код редне везе еквивалентна отпорност је увек \textbf{већа од највеће}
          појединачне отпорности.
    \item Ако је редно везано $n$ једнаких отпорника отпорности $R$, онда је
          $R_e = n \cdot R$.
\end{itemize}

\begin{primer}
Три отпорника, $R_1 = 10\,\Omega$, $R_2 = 20\,\Omega$ и $R_3 = 30\,\Omega$,
везана су редно. Израчунати еквивалентну отпорност.

\textbf{Решење:}
\[
R_e = R_1 + R_2 + R_3 = 10\,\Omega + 20\,\Omega + 30\,\Omega = 60\,\Omega
\]
Резултат је већи од највеће отпорности ($30\,\Omega$), као што и треба да буде.
\end{primer}

\begin{primer}
Отпорници $R_1 = 1{,}5\,\text{k}\Omega$ и $R_2 = 500\,\Omega$ везани су редно.
Израчунати еквивалентну отпорност.

\textbf{Решење:}

Пре сабирања све отпорности морају бити у истој јединици:
$R_1 = 1{,}5\,\text{k}\Omega = 1500\,\Omega$.
\[
R_e = R_1 + R_2 = 1500\,\Omega + 500\,\Omega = 2000\,\Omega = 2\,\text{k}\Omega
\]
\end{primer}

%==================================================================
\section{Паралелна веза отпорника}
%==================================================================

Отпорници су везани \textbf{паралелно} када су сви прикључени између
\textbf{истих тачака} $A$ и $B$. Струја тада има више путева: у тачки $A$ се
грана, а у тачки $B$ поново спаја.

\begin{center}
\begin{circuitikz}
    \draw (0,1) node[left]{$A$} to[short, o-*] (1,1);
    \draw (1,1) -- (1,2) to[R, l=$R_1$] (4,2) -- (4,1);
    \draw (1,1) -- (1,0) to[R, l=$R_2$] (4,0) -- (4,1);
    \draw (4,1) to[short, *-o] (5,1) node[right]{$B$};
\end{circuitikz}
\end{center}

Сви паралелни отпорници су између истих тачака, па је \textbf{на свима исти
напон}. Сваки нови отпорник отвара још један пут за струју, па укупна
отпорност \textbf{опада}. Овде се не сабирају отпорности, већ њихове
реципрочне вредности:

\begin{kljucno}
\[
\frac{1}{R_e} = \frac{1}{R_1} + \frac{1}{R_2} + \dots + \frac{1}{R_n}
\]
\end{kljucno}

За \textbf{два} отпорника формула се може написати лакше („производ кроз збир”):
\[
R_e = \frac{R_1 \cdot R_2}{R_1 + R_2}
\]

\begin{itemize}
    \item Код паралелне везе еквивалентна отпорност је увек \textbf{мања од
          најмање} појединачне отпорности.
    \item Ако је паралелно везано $n$ једнаких отпорника отпорности $R$, онда
          је $R_e = \dfrac{R}{n}$.
\end{itemize}

\begin{primer}
Отпорници $R_1 = 12\,\Omega$ и $R_2 = 4\,\Omega$ везани су паралелно.
Израчунати еквивалентну отпорност.

\textbf{Решење:}
\[
R_e = \frac{R_1 \cdot R_2}{R_1 + R_2}
    = \frac{12\,\Omega \cdot 4\,\Omega}{12\,\Omega + 4\,\Omega}
    = \frac{48}{16}\,\Omega = 3\,\Omega
\]
Резултат је мањи од најмање отпорности ($4\,\Omega$), као што и треба да буде.
\end{primer}

\begin{primer}
Три отпорника, $R_1 = 6\,\Omega$, $R_2 = 3\,\Omega$ и $R_3 = 2\,\Omega$, везана
су паралелно. Израчунати еквивалентну отпорност.

\begin{center}
\begin{circuitikz}
    \draw (0,1) node[left]{$A$} to[short, o-*] (1,1);
    \draw (1,1) -- (1,2) to[R, l=$R_1$] (4,2) -- (4,1);
    \draw (1,1) to[R, l=$R_2$, -*] (4,1);
    \draw (1,1) -- (1,0) to[R, l=$R_3$] (4,0) -- (4,1);
    \draw (4,1) to[short, -o] (5,1) node[right]{$B$};
\end{circuitikz}
\end{center}

\textbf{Решење:}

Код три отпорника користимо основну формулу и сводимо разломке на заједнички
именилац:
\[
\frac{1}{R_e} = \frac{1}{6\,\Omega} + \frac{1}{3\,\Omega} + \frac{1}{2\,\Omega}
              = \frac{1}{6\,\Omega} + \frac{2}{6\,\Omega} + \frac{3}{6\,\Omega}
              = \frac{6}{6\,\Omega} = \frac{1}{1\,\Omega}
\]
Добили смо $\dfrac{1}{R_e}$, па резултат још треба окренути:
\[
R_e = 1\,\Omega
\]
\end{primer}

\begin{primer}
Три једнака отпорника, сваки од $30\,\Omega$, везана су паралелно.
Израчунати еквивалентну отпорност.

\textbf{Решење:}
\[
R_e = \frac{R}{n} = \frac{30\,\Omega}{3} = 10\,\Omega
\]
\end{primer}

%==================================================================
\section{Везе кондензатора}
%==================================================================

Кондензатори се, као и отпорници, могу везати паралелно и редно. Важно је
запамтити да су формуле \textbf{обрнуте} у односу на отпорнике: оно што се код
отпорника рачуна за редну везу, код кондензатора важи за паралелну, и обрнуто.

\subsection{Паралелна веза кондензатора}

\begin{center}
\begin{circuitikz}
    \draw (0,1) node[left]{$A$} to[short, o-*] (1,1);
    \draw (1,1) -- (1,2) to[C, l=$C_1$] (4,2) -- (4,1);
    \draw (1,1) -- (1,0) to[C, l=$C_2$] (4,0) -- (4,1);
    \draw (4,1) to[short, *-o] (5,1) node[right]{$B$};
\end{circuitikz}
\end{center}

Сви кондензатори су између истих тачака, па је \textbf{на свима исти напон}
$U$. Сваки кондензатор прима своје наелектрисање, а укупно наелектрисање је
њихов збир: $Q = Q_1 + Q_2 + \dots + Q_n$. Зато се капацитивности сабирају:

\begin{kljucno}
\[
C_e = C_1 + C_2 + \dots + C_n
\]
\end{kljucno}

Паралелно везани кондензатори понашају се као један кондензатор са
\textbf{већом површином плоча}, а већа површина значи већу капацитивност.
Еквивалентна капацитивност је зато \textbf{већа од највеће} појединачне.

\begin{primer}
Кондензатори $C_1 = 2\,\mu\text{F}$ и $C_2 = 3\,\mu\text{F}$ везани су
паралелно. Израчунати еквивалентну капацитивност.

\textbf{Решење:}
\[
C_e = C_1 + C_2 = 2\,\mu\text{F} + 3\,\mu\text{F} = 5\,\mu\text{F}
\]
\end{primer}

\subsection{Редна веза кондензатора}

\begin{center}
\begin{circuitikz}
    \draw (0,0) node[left]{$A$} to[short, o-] (0.5,0)
          to[C, l=$C_1$] (2.5,0)
          to[C, l=$C_2$] (4.5,0)
          to[short, -o] (5,0) node[right]{$B$};
\end{circuitikz}
\end{center}

Код редне везе \textbf{сви кондензатори имају исто наелектрисање} $Q$, а
напони на њима се сабирају: $U = U_1 + U_2 + \dots + U_n$. Капацитивности се
овде рачунају као отпорности у паралелној вези:

\begin{kljucno}
\[
\frac{1}{C_e} = \frac{1}{C_1} + \frac{1}{C_2} + \dots + \frac{1}{C_n}
\]
\end{kljucno}

За \textbf{два} кондензатора:
\[
C_e = \frac{C_1 \cdot C_2}{C_1 + C_2}
\]

Редно везани кондензатори понашају се као један кондензатор са \textbf{већим
растојањем између плоча}, а веће растојање значи мању капацитивност.
Еквивалентна капацитивност је зато \textbf{мања од најмање} појединачне.

\begin{primer}
Кондензатори $C_1 = 6\,\mu\text{F}$ и $C_2 = 3\,\mu\text{F}$ везани су редно.
Израчунати еквивалентну капацитивност.

\textbf{Решење:}
\[
C_e = \frac{C_1 \cdot C_2}{C_1 + C_2}
    = \frac{6\,\mu\text{F} \cdot 3\,\mu\text{F}}{6\,\mu\text{F} + 3\,\mu\text{F}}
    = \frac{18}{9}\,\mu\text{F} = 2\,\mu\text{F}
\]
\end{primer}



%==================================================================
\section{Резиме}
%==================================================================

\begin{center}
\renewcommand{\arraystretch}{2}
\begin{tabular}{lcc}
    \toprule
                          & \textbf{Отпорници}
                          & \textbf{Кондензатори} \\
    \midrule
    Редна веза            & $R_e = R_1 + R_2$
                          & $\dfrac{1}{C_e} = \dfrac{1}{C_1} + \dfrac{1}{C_2}$ \\
    Паралелна веза        & $\dfrac{1}{R_e} = \dfrac{1}{R_1} + \dfrac{1}{R_2}$
                          & $C_e = C_1 + C_2$ \\
    Исто у редној вези    & струја $I$     & наелектрисање $Q$ \\
    Исто у паралелној вези & напон $U$     & напон $U$ \\
    \bottomrule
\end{tabular}
\end{center}

\end{document}